% ============================================================
%  Chapter 4 — Case Study: M/V Obbola and the Baltic Sea Loop
% ============================================================
\section{Case Study: M/V Obbola and the Baltic Sea Loop}
\label{ch:casestudy}

% ── 4.1 Vessel description ────────────────────────────────────
\subsection{Vessel Description}
\label{sec:case:vessel}

M/V~Obbola is a \RoRo{} vessel operated by SCA Logistics on a cyclic
Baltic Sea service. Key technical parameters are summarised in
\cref{tab:vessel_params}.

\begin{table}[h]
\centering
\caption{M/V Obbola — principal technical parameters}
\label{tab:vessel_params}
\begin{tabular}{lrl}
  \toprule
  Parameter & Value & Unit \\
  \midrule
  IMO number            & 9\,395\,343 & --- \\
  Vessel type           & \RoRo{}     & --- \\
  Gross Tonnage (\GT{}) & 19\,918     & GT \\
  Deadweight Tonnage    & 11\,446     & t \\
  Cargo capacity        & 9\,730      & t \\
  Length (L\textsubscript{pp}) & 160.4 & m \\
  Breadth               & 23.5        & m \\
  Summer draught        & 6.7         & m \\
  Ice class             & 1A          & --- \\
  $f_m$ correction      & 1.05        & --- \\
  $f_i$ correction      & 1.00        & --- \\
  Effective capacity ($f_i \cdot f_m \cdot \GT{}$)
                        & 20\,913.9   & GT \\
  Design speed (summer, light) & 16.0 & kn \\
  Fuel type             & HFO         & --- \\
  CO$_2$ emission factor $\gamma$ & 3\,114 & kg/t\textsubscript{fuel} \\
  \bottomrule
\end{tabular}
\end{table}

% ── 4.2 CII parameterisation ──────────────────────────────────
\subsection{CII Parameterisation}
\label{sec:case:cii}

The required \CII{} for M/V~Obbola is computed from the \IMO{}
reference formula for \RoRo{} vessels:

\begin{equation}
  \ciireq^{(Y)} = a \cdot \GT^{-c} \cdot (1 - r^{(Y)})
  \label{eq:cii_required}
\end{equation}

\noindent with $a = 1967$, $c = 0.485$ \parencite{mepc336}, giving a
reference value $\ciireq^{\mathrm{ref}} = 16.17$~\gGTnm{}.
The year-specific required \CII{} values after applying the annual
reduction factors $r^{(Y)}$ are given in \cref{tab:cii_required}.

\begin{table}[h]
\centering
\caption{Required CII and scenario thresholds for M/V Obbola}
\label{tab:cii_required}
\begin{tabular}{lrrrr}
  \toprule
  Year & {Reduction (\%)} & {Required CII} & {B boundary ($\times 0.89$)}
       & {C boundary ($\times 1.08$)} \\
  \midrule
  2023 & 5.0  & 15.3601 & 13.6705 & 16.5889 \\
  2024 & 7.0  & 15.0367 & 13.3827 & 16.2397 \\
  2025 & 9.0  & 14.7134 & 13.0949 & 15.8904 \\
  2026 & 11.0 & 14.3900 & 12.8071 & 15.5412 \\
  2030 & 31.0 & 11.1563 & 9.9291  & 12.0488 \\
  \bottomrule
\end{tabular}
\smallskip\\
{\footnotesize Rating boundaries use RoRo-specific multipliers from
MEPC.339(76): $d_1=0.76$, $d_2=0.89$, $d_3=1.08$, $d_4=1.27$.}
\end{table}

% ── 4.3 Route and network ─────────────────────────────────────
\subsection{Baltic Sea Loop}
\label{sec:case:route}

M/V~Obbola operates a cyclic Baltic Sea service on behalf of SCA, calling
at Swedish forestry and paper-industry ports before transiting to northern
European hub ports.  The service connects SCA's Swedish production sites ---
Umeå (SEUME), Sundsvall (SETUN / SESDL), Iggesund (SEIGG), Mälarhamn
(SEMMA), Helsingborg (SEHLD / SEHEL), and Oxelösund (SEOXE) --- with
continental and UK ports at Kiel (DEKEL), Rotterdam (NLRTM), and Tilbury
(GBTIL).  The loop is completed approximately every three weeks
(observed mean cycle time approximately 935~h across arcs), corresponding
to roughly 17 complete loops per calendar year.

The EU \MRV{} dataset for December 2021 -- August 2022 contains 153
individual voyage legs spanning 25 unique arc types (origin--destination
port pairs) before filtering, of which 22 pass the minimum observation
threshold ($n \geq 2$).  The total observed distance across these arcs
is 26{,}575~nm, and the per-loop distance used in the model
(the expected total distance for one complete service cycle)
is approximately 10{,}540~nm.  \Cref{fig:route_map} illustrates the
network topology.

\begin{figure}[h]
  \centering
  % Insert: route map of the Baltic Sea loop (PNG/PDF)
  % Placeholder — replace \fbox with \includegraphics{Graphics/route_map}
  % once the figure is generated.
  \fbox{\parbox{0.85\textwidth}{\centering\vspace{2cm}
        \textit{Figure: Baltic Sea loop --- M/V Obbola service network.}\\
        \smallskip
        Ports: SEUME (Umeå), SETUN/SESDL (Sundsvall), SEIGG (Iggesund),\\
        SEMMA (Mälarhamn), SEHLD/SEHEL (Helsingborg), SEOXE (Oxelösund),\\
        DEKEL (Kiel), NLRTM (Rotterdam), GBTIL (Tilbury).
        \vspace{2cm}}}
  \caption{Baltic Sea loop operated by M/V~Obbola (SCA Logistics).
           Solid arrows indicate the primary service arcs; dashed arrows
           indicate alternative routing options observed in the EU \MRV{} data.}
  \label{fig:route_map}
\end{figure}

% ── 4.4 EU MRV data ───────────────────────────────────────────
\subsection{EU MRV Operational Data}
\label{sec:case:mrv}

EU \MRV{} voyage data for M/V~Obbola covering December 2021 through
August 2022 were used to calibrate the network and validate the model.
The dataset comprises 153 individual voyage legs across 22 unique arc
types (after applying a minimum observation filter of $n \geq 2$).

\begin{table}[h]
\centering
\caption{EU MRV dataset summary}
\label{tab:mrv_summary}
\begin{tabular}{lr}
  \toprule
  Item & Value \\
  \midrule
  Reporting period & Dec 2021 -- Aug 2022 \\
  Total voyage legs & 153 \\
  Unique arc types ($n \geq 2$) & 22 \\
  Total distance sailed & 26\,575 nm \\
  Total CO$_2$ reported & 6\,352.8 t \\
  Attained CII (2022, reported) & 11.43 g/GT$\cdot$nm \\
  IMO CII rating (2022) & B \\
  Mean cargo per leg & 50.8\% capacity \\
  \bottomrule
\end{tabular}
\end{table}

% ── 4.5 Schedule and time budgets ─────────────────────────────
\subsection{Schedule and Time Budgets}
\label{sec:case:schedule}

Per-arc time budgets $\budget$ are derived from the EU \MRV{} data as:

\begin{equation}
  \budget = \bar{t}^{\mathrm{transit}}_{ij}
            + \bar{t}^{\mathrm{wait}}_{ij}
  \label{eq:budget_def}
\end{equation}

\noindent where $\bar{t}^{\mathrm{transit}}$ is the mean observed transit
time and $\bar{t}^{\mathrm{wait}}$ is the mean observed port waiting time
at the arrival port. This budget represents the maximum time available
before the vessel must depart the next port on the published schedule.

Three arcs have particularly tight budgets that force high minimum speeds:

\begin{table}[h]
\centering
\caption{Binding schedule constraints — bottleneck arcs}
\label{tab:bottleneck_arcs}
\begin{tabular}{lrrrr}
  \toprule
  Arc & {Distance (nm)} & {Budget (h)} & {Min speed (kn, continuous)}
      & {Min speed (kn, grid)} \\
  \midrule
  SEMMA $\to$ SEHLD & 692 & 48.0 & 14.41 & 15 \\
  SEMMA $\to$ SEUME & 690 & 48.0 & 14.37 & 15 \\
  SEIGG $\to$ DEKEL & 659 & 48.0 & 13.74 & 14 \\
  SETUN $\to$ NLRTM & 1047 & 76.5 & 13.69 & 14 \\
  \bottomrule
\end{tabular}
\smallskip\\
{\footnotesize These arcs correspond to published 48-hour overnight
schedule windows. They define the minimum achievable CII under the
fixed schedule.}
\end{table}

% ── 4.6 Model validation ──────────────────────────────────────
\subsection{Model Validation}
\label{sec:case:validation}

The model is validated in Scenario~S1 (as-is replay) by comparing the
attained \CII{} predicted by the model against the auditor-reported
value from the EU \MRV{} dataset. Using mean observed speeds on each arc,
the model predicts an attained \CII{} of 11.29~\gGTnm{}, compared to
the reported value of 11.43~\gGTnm{} --- a relative error of 1.2\%.
Both values correspond to a Rating~B under the RoRo-specific boundary
multipliers from MEPC.339(76). This close agreement validates the
fuel consumption model and network parameterisation.

The three sources of discrepancy enumerated above (speed rounding, per-arc
mean cargo, exclusion of port auxiliary consumption) are all conservative in
the same direction: they collectively produce a slight underestimate of
emissions.  A 1.2\,\% relative error is well within the range considered
acceptable for operational planning models calibrated to aggregate MRV data;
\textcite{psaraftis2013speed} note that fuel model errors of 5--15\,\% are
common in fleet-level analyses, making sub-2\,\% accuracy at the route level
highly satisfactory.

The dominant contributor to the residual discrepancy is the exclusion of
port-side auxiliary engine consumption.  Auxiliary power draw during loading
and unloading operations is recorded in the MRV dataset under the port call
emission category but is not captured by the underway fuel model, which covers
only propulsion fuel.  Were auxiliary consumption included, the model's
predicted \CII{} would move closer to the reported value.  Given that this
thesis focuses on the operational levers available for underway speed
optimisation, the exclusion is methodologically appropriate.
